% GENERATED FILE. Edit canonical lessons, then run npm run generate:books.
% canonical-source-hash: fnv1a64:09e61ced5f6a8a3f
% canonical-lessons: HI-C324-praja, HI-C324-hukm, HI-C324-gulam, HI-C324-vidroh, HI-C324-lalkarna

\chapter{Rebellion}
\label{ch:hi-rebellion}
\hlchaptermodality{324}

\begin{chapteropening}
\textbf{By the end of this chapter:} \emph{I can tell how people rose against a ruler.}

Complete the last lesson of chapter 324: I can tell how people rose against a ruler.
\end{chapteropening}

\section[\texorpdfstring{\hi{प्रजा} (prajā)}{प्रजा (prajā)}]{\hi{प्रजा} (prajā) — the subjects (of a king)}

\label{lesson:HI-C324-praja}



\begin{quote}
Before the new one: say the Hindi for a martyr, then the Hindi for a treaty.
\end{quote}

\subsection*{You'll want to know: \hi{प्रजा}}
\textbf{\hi{प्रजा}} — \emph{prajā} — \textquotedblleft{}the subjects (of a king)\textquotedblright{}.

\begin{quote}
\hi{राजा} \hi{की} \hi{प्रजा} \hi{बहुत} \hi{ग़रीब} \hi{थी।}
\end{quote}

\begin{quote}
\emph{rājā kī prajā bahut ġarīb thī.}
\end{quote}

\begin{quote}
The king's subjects were very poor.
\end{quote}

Feminine.

\begin{cousinweb}
\textbf{\hi{प्रजा}} (\emph{prajā}) is Sanskrit \emph{prajā}, \textquotedblleft{}offspring, people, subjects\textquotedblright{} (Monier-Williams): the people a king rules. It is singular and feminine: \textbf{\hi{प्रजा} \hi{ग़रीब} \hi{थी}}.
\end{cousinweb}

\begin{grammarlens}[title={What you've built}]
One more word for telling what happened.
\end{grammarlens}

\subsection*{Your turn}
\begin{itemize}
\raggedright
  \item \emph{Say it:} \emph{prajā}
  \item \emph{Say it:} \emph{prajā}, once more
  \item \emph{Say it:} \emph{rājā kī prajā bahut ġarīb thī.} — The king's subjects were very poor.
  \item \emph{Recall:} say the Hindi for a martyr, then the Hindi for a treaty, then say \emph{prajā} again
\end{itemize}

\begin{tcolorbox}[breakable,colback=teal!4,colframe=teal!35!black,title={Before you move on}]
What does \hi{प्रजा} mean? (\textquotedblleft{}The subjects (of a king)\textquotedblright{}.) Say it once more.
\end{tcolorbox}
\section[\texorpdfstring{\hi{हुक्म} (hukm)}{हुक्म (hukm)}]{\hi{हुक्म} (hukm) — an order; a command}

\label{lesson:HI-C324-hukm}



\begin{quote}
Before the new one: say the Hindi for a treaty, then the Hindi for the subjects (of a king).
\end{quote}

\subsection*{You'll want to know: \hi{हुक्म}}
\textbf{\hi{हुक्म}} — \emph{hukm} — \textquotedblleft{}an order; a command\textquotedblright{}.

\begin{quote}
\hi{राजा} \hi{के} \hi{हुक्म} \hi{से} \hi{सिपाही} \hi{गाँव} \hi{में} \hi{आए।}
\end{quote}

\begin{quote}
\emph{rājā ke hukm se sipāhī gā̃v meṁ āe.}
\end{quote}

\begin{quote}
On the king's orders, soldiers came to the village.
\end{quote}

Masculine.

\begin{cousinweb}
\textbf{\hi{हुक्म}} (\emph{hukm}) is usually traced, through Persian, to Arabic \emph{ḥukm}, \textquotedblleft{}a judgement, a ruling\textquotedblright{}. \textbf{\hi{हुक्म} \hi{देना}} is to give an order.
\end{cousinweb}

\begin{grammarlens}[title={What you've built}]
One more word for telling what happened.
\end{grammarlens}

\subsection*{Your turn}
\begin{itemize}
\raggedright
  \item \emph{Say it:} \emph{hukm}
  \item \emph{Say it:} \emph{hukm}, once more
  \item \emph{Say it:} \emph{rājā ke hukm se sipāhī gā̃v meṁ āe.} — On the king's orders, soldiers came to the village.
  \item \emph{Recall:} say the Hindi for a treaty, then the Hindi for the subjects (of a king), then say \emph{hukm} again
\end{itemize}

\begin{tcolorbox}[breakable,colback=teal!4,colframe=teal!35!black,title={Before you move on}]
What does \hi{हुक्म} mean? (\textquotedblleft{}An order; a command\textquotedblright{}.) Say it once more.
\end{tcolorbox}
\section[\texorpdfstring{\hi{ग़ुलाम} (ġulām)}{ग़ुलाम (ġulām)}]{\hi{ग़ुलाम} (ġulām) — a slave}

\label{lesson:HI-C324-gulam}



\begin{quote}
Before the new one: say the Hindi for the subjects (of a king), then the Hindi for an order.
\end{quote}

\subsection*{You'll want to know: \hi{ग़ुलाम}}
\textbf{\hi{ग़ुलाम}} — \emph{ġulām} — \textquotedblleft{}a slave\textquotedblright{}.

\begin{quote}
\hi{उस} \hi{राजा} \hi{के} \hi{महल} \hi{में} \hi{सौ} \hi{ग़ुलाम} \hi{थे।}
\end{quote}

\begin{quote}
\emph{us rājā ke mahal meṁ sau ġulām the.}
\end{quote}

\begin{quote}
There were a hundred slaves in that king's palace.
\end{quote}

Masculine.

\begin{cousinweb}
\textbf{\hi{ग़ुलाम}} (\emph{ġulām}) is usually traced, through Persian, to Arabic \emph{ghulām}, \textquotedblleft{}a youth, a servant\textquotedblright{}. In the history of freedom movements it also means a people ruled by others.
\end{cousinweb}

\begin{grammarlens}[title={What you've built}]
One more word for telling what happened.
\end{grammarlens}

\subsection*{Your turn}
\begin{itemize}
\raggedright
  \item \emph{Say it:} \emph{ġulām}
  \item \emph{Say it:} \emph{ġulām}, once more
  \item \emph{Say it:} \emph{us rājā ke mahal meṁ sau ġulām the.} — There were a hundred slaves in that king's palace.
  \item \emph{Recall:} say the Hindi for the subjects (of a king), then the Hindi for an order, then say \emph{ġulām} again
\end{itemize}

\begin{tcolorbox}[breakable,colback=teal!4,colframe=teal!35!black,title={Before you move on}]
What does \hi{ग़ुलाम} mean? (\textquotedblleft{}A slave\textquotedblright{}.) Say it once more.
\end{tcolorbox}
\section[\texorpdfstring{\hi{विद्रोह} (vidroh)}{विद्रोह (vidroh)}]{\hi{विद्रोह} (vidroh) — a rebellion; a revolt}

\label{lesson:HI-C324-vidroh}



\begin{quote}
Before the new one: say the Hindi for an order, then the Hindi for a slave.
\end{quote}

\subsection*{You'll want to know: \hi{विद्रोह}}
\textbf{\hi{विद्रोह}} — \emph{vidroh} — \textquotedblleft{}a rebellion; a revolt\textquotedblright{}.

\begin{quote}
\hi{भूख} \hi{की} \hi{वजह} \hi{से} \hi{किसानों} \hi{ने} \hi{विद्रोह} \hi{किया।}
\end{quote}

\begin{quote}
\emph{bhūkh kī vajah se kisānoṁ ne vidroh kiyā.}
\end{quote}

\begin{quote}
The farmers rose up because they were going hungry.
\end{quote}

Masculine.

\begin{cousinweb}
\textbf{\hi{विद्रोह}} (\emph{vidroh}) is built from \emph{vi-} and Sanskrit \emph{droha}, \textquotedblleft{}injury, treachery\textquotedblright{}; Monier-Williams gives \emph{droha} and the verb \emph{vi-druh}, \textquotedblleft{}to injure, to do wrong\textquotedblright{}, but not \emph{vidroha} itself, and the sense of a revolt is probably the later language's own. Hindi uses it for a revolt: \textbf{\hi{विद्रोह} \hi{करना}} is to rise up against a ruler.
\end{cousinweb}

\begin{grammarlens}[title={What you've built}]
One more word for telling what happened.
\end{grammarlens}

\subsection*{Your turn}
\begin{itemize}
\raggedright
  \item \emph{Say it:} \emph{vidroh}
  \item \emph{Say it:} \emph{vidroh}, once more
  \item \emph{Say it:} \emph{bhūkh kī vajah se kisānoṁ ne vidroh kiyā.} — The farmers rose up because they were going hungry.
  \item \emph{Recall:} say the Hindi for an order, then the Hindi for a slave, then say \emph{vidroh} again
\end{itemize}

\begin{tcolorbox}[breakable,colback=teal!4,colframe=teal!35!black,title={Before you move on}]
What does \hi{विद्रोह} mean? (\textquotedblleft{}A rebellion; a revolt\textquotedblright{}.) Say it once more.
\end{tcolorbox}
\section[\texorpdfstring{\hi{ललकारना} (lalkārnā)}{ललकारना (lalkārnā)}]{\hi{ललकारना} (lalkārnā) — to challenge; to call out (an enemy)}

\label{lesson:HI-C324-lalkarna}



\begin{quote}
Before the new one: say the Hindi for a slave, then the Hindi for a rebellion.
\end{quote}

\subsection*{You'll want to know: \hi{ललकारना}}
\textbf{\hi{ललकारना}} — \emph{lalkārnā} — \textquotedblleft{}to challenge; to call out (an enemy)\textquotedblright{}.

\begin{quote}
\hi{एक} \hi{किसान} \hi{ने} \hi{राजा} \hi{के} \hi{सिपाहियों} \hi{को} \hi{ललकारा।}
\end{quote}

\begin{quote}
\emph{ek kisān ne rājā ke sipāhiyoṁ ko lalkārā.}
\end{quote}

\begin{quote}
A farmer challenged the king's soldiers.
\end{quote}

Past, with \hi{ने} (he or she …): \hi{उसने} \hi{दुश्मन} \hi{को} \hi{ललकारा}, \emph{usne duśman ko lalkārā}.

\begin{cousinweb}
\textbf{\hi{ललकारना}} (\emph{lalkārnā}) is the verb of \emph{lalkār}, a challenge, which CDIAL 10973 lists under \emph{*lallakka}, a noise made to threaten: to shout at an enemy to come and fight.
\end{cousinweb}

\begin{grammarlens}[title={What you've built}]
One more word for telling what happened.
\end{grammarlens}

\subsection*{Your turn}
\begin{itemize}
\raggedright
  \item \emph{Say it:} \emph{lalkārnā}
  \item \emph{Say it:} \emph{lalkārnā}, once more
  \item \emph{Say it:} \emph{ek kisān ne rājā ke sipāhiyoṁ ko lalkārā.} — A farmer challenged the king's soldiers.
  \item \emph{Recall:} say the Hindi for a slave, then the Hindi for a rebellion, then say \emph{lalkārnā} again
\end{itemize}

\begin{tcolorbox}[breakable,colback=teal!4,colframe=teal!35!black,title={Before you move on}]
What does \hi{ललकारना} mean? (\textquotedblleft{}To challenge; to call out (an enemy)\textquotedblright{}.) Say it once more.
\end{tcolorbox}
